\section{モデルの重みを直接操作する研究}
MIT CSAILの研究者による「Weight-Space Programming」の紹介投稿が広く拡散された。Grokが記録したブログの実験では、GPT-6 Astraが別のニューラルネットワークの重みを読み書きし、Qwen3.5-0.8Bの重みの0.016\%を変更することで、MATH-500の成績が7.9\%から48.1\%になったと報告されている。ゲーム用ポリシーの識別や失敗修正、重みの新規生成も紹介された。従来のプロンプトや追加学習とは異なるモデル操作として注目される一方、正式な論文公開や第三者による再現はGrokの観測範囲で未確認である。
\dailyxsource{Rikiya Takehi、11月3日02:51 JST}{https://x.com/rikiyatakehi/status/2117313742891327994}

\section{Mistral Large 4の技術レポート}
Mistralの研究責任者による投稿で、Mistral Large 4（ML4、通称Le Chonk）の技術レポートが共有された。Grokによれば、モデルのオープンウェイト公開は10月末で今回の観測窓外、技術レポートの公開は11月2日である。投稿では、プレビュー版より効率的かつコンパクトになり、コーディングエージェントやサイバー領域で高い性能を示すと説明された。性能面の主張は提供側の説明であり、独立ベンチマークによる比較は今後の確認事項となる。
\dailyxsource{Pierre Stock、11月3日01:30 JST}{https://x.com/PierreStock/status/2117293394238013698}

\section{次世代GPTの公開時期を巡る発言}
ADIPECでSam Altmanが、Navier--Stokes方程式の証明に使われたとされる内部モデルについて、次世代GPTとして近く公開する旨を述べたという投稿が拡散した。投稿には、1万のエージェントが4日間協働したとの引用も含まれる。ただし、Grokは公式発表や一次映像を直接確認しておらず、モデルの同一性、証明の妥当性、具体的な公開日は未確認である。新モデルの発売予告として断定せず、X上で注目された発言の紹介として扱う。
\dailyxsource{prinz、11月2日22:59 JST}{https://x.com/deredleritt3r/status/2117255388353663125}

\section{Beaconによるエージェントの記憶共有}
Asymptote Labsが、コーディングエージェント向けのオープンソースメモリツール「Beacon」を発表した。Claude CodeやCursorなどで行った過去の作業を再利用し、セッションで得た知見をスキルとして蓄積する構想である。開発者は、繰り返し説明するトークン消費を抑え、エージェントが過去の経験から改善できると主張する。GrokはGitHubリポジトリへの言及を確認したが、実際の削減効果や運用上の安定性までは精査していない。
\dailyxsource{Justin D'Souza、11月3日00:59 JST}{https://x.com/jqdsouza/status/2117285537836908696}

\section{GitHub Spec Kitと仕様駆動開発}
GitHubの「Spec Kit」を紹介する投稿が、曖昧な指示だけでコードを生成する手法に代わる仕様駆動ワークフローとして話題となった。実装前に要求や設計を明文化し、AIコーディングエージェントの作業を構造化する狙いがある。投稿ではGitHubのスター数が13万を超え、38の統合に対応すると主張されているが、Grokは公式リポジトリの数値や対応範囲を直接確認していない。仕様の品質と生成コードの整合性をどう保つかが注目点となる。
\dailyxsource{Jainam Parmar、11月2日22:12 JST}{https://x.com/aiwithjainam/status/2117243550098424270}

\section{DeepSeek HarnessのCreator Mode}
DeepSeek Harnessに、会話からプラグインを生成してインストールできる「Creator Mode」が追加されたとする投稿が共有された。公式アカウントのデモ動画が引用され、機能をプラグインとして拡張していく設計が紹介されている。開発者側は「Everything is a plugin」という方向性を掲げるが、Grokは詳細なドキュメントや動作の安定性を確認していない。実装の自由度に加え、生成プラグインの権限管理や再現性が今後の検討事項となる。
\dailyxsource{DeepSeek Harness、観測窓内の引用元投稿}{https://x.com/DeepSeekHarness/status/2117256196227019194}

\section{AI学習オプトアウト信号の実装時期}
robots.txtにおけるAI学習用途の統一的な意思表示を巡り、\texttt{\_\_usage\_ai-training} トークンの実装時期が議論された。Grokの収集では、Google、Microsoft、OpenAI、Anthropicが対応を表明した10月30日の発表は観測窓外で、今回の対象は11月3日朝の反応投稿である。Anthropicが先行し、他社は2027年になるとの報道が紹介され、共通仕様への賛同と実装速度の差が論点となった。各社の実際の対応状況は継続確認が必要である。
\dailyxsource{SiMelt、11月3日06:59 JST}{https://x.com/zshjzls/status/2117376197470662979}

\section{OpenCodeのマルチモデル対応を紹介}
オープンソースのコーディングエージェント「OpenCode」が、複数のモデルやプロバイダを切り替えて利用できる選択肢として紹介された。投稿では75以上のプロバイダ対応、MITライセンス、既存サブスクリプションの利用、コードを保存しないプライバシー重視の設計が挙げられた。ただし、これらはGrokが収集した紹介投稿の主張であり、対応数やデータ取扱いの詳細は未精査である。利用者がモデルと実行環境を選べることへの関心がうかがえる。
\dailyxsource{Akshay Shinde、11月2日21:30 JST}{https://x.com/ConsciousRide/status/2117232811858145538}

\section{TermiX Skillsのマーケットプレイス}
TermiXの「Agent Skills v1.9.0」により、既存エージェントをマーケットプレイスへ登録し、外部からタスクと報酬を受け取れるようになったとする投稿が共有された。単なる個人用アシスタントから、仕事を受託するワーカーへの展開を示唆する内容である。ただし、Grokは公式な機能仕様や実際の取引実績を確認していない。エージェントの公開・決済・実行を結び付ける試みとして注目されるが、現時点では未確認情報として扱う。
\dailyxsource{Ethreums、11月2日13:59 JST}{https://x.com/HVnS42442600/status/2117119370363392439}

\section{WonderDriveによる共有コンテキスト}
WonderSearchAIが、複数のAIエージェント向け共有ストレージ「WonderDrive」を発表したとの投稿があった。ClaudeとCodexの間で作業コンテキストを共有し、異なるエージェント間の引き継ぎを容易にするという。投稿者は、埋め込みモデルを使わずに情報を照合できると主張しているが、技術的な仕組みや性能はGrokの観測では未確認である。エージェントを横断した記憶とデータ共有への需要を示す事例として記録する。
\dailyxsource{dale、11月2日12:01 JST}{https://x.com/daleverett/status/2117089766084882854}
